۹

«لِی‌گاید»، خلاصه‌ی «راهنمای چگونگی خواباندن دختران». در ۱۵۰ صفحه‌ی سوزان، او گفت که خرد جمعی ده‌ها هنرمند جلب توجه، که برای نزدیک به یک دهه دانش خود را در گروه‌های خبری تبادل می‌کردند، مخفیانه در حال کار بودند تا هنر اغواگری را به یک علم دقیق تبدیل کنند. اطلاعات نیاز به بازنویسی و سازماندهی به یک کتاب چگونگی منسجم داشت و او فکر می‌کرد من فردی مناسب برای انجام این کار هستم.

من مطمئن نبودم. می‌خواهم ادبیات بنویسم، نه اینکه به نوجوانان هوسران مشاوره‌ بدهم. اما، البته، به او گفتم نگاه کردن به آن ضرری ندارد.

لحظه‌ای که شروع به خواندن کردم، زندگی‌ام تغییر کرد. بیشتر از هر کتاب یا سند دیگری — چه کتاب مقدس، چه جنایت و مکافات، یا لذت آشپزی — «لِی‌گاید» چشمانم را باز کرد. و نه لزوماً به‌خاطر اطلاعات درون آن، بلکه به‌خاطر راهی که مرا به پایین کشاند.

وقتی به سال‌های نوجوانی‌ام نگاه می‌کنم، یک افسوس عمده دارم و آن هم ارتباطی با درس نخواندن به‌اندازه‌ی کافی، مهربان نبودن با مادرم، یا کوبیدن ماشین پدرم به یک اتوبوس عمومی ندارد. فقط این که با آن‌قدر دختر بازی نکرده بودم. من مرد عمیقی هستم — هر سه سال یک‌بار برای سرگرمی، کتاب «اولیس جیمز جویس» را مجدداً مطالعه می‌کنم. خودم را به اندازه‌ی معقولی شهودی می‌دانم. در اصل آدم خوبی هستم و سعی می‌کنم به دیگران آسیب نرسانم. اما به نظر می‌رسد نمی‌توانم به حالت بالاتر وجودی پیشرفت کنم چون زمان بسیار زیادی را به فکر کردن درباره‌ی زنان صرف می‌کنم.

و می‌دانم که تنها نیستم. وقتی اولین بار هیو هفنر را ملاقات کردم، او هفتاد و سه ساله بود. به گفته‌ی خودش با بیش از هزار نفر از زیباترین زنان جهان خوابیده بود، اما همه‌ی چیزی که می‌خواست درباره‌اش صحبت کند، سه دوست دخترش—مندی، برند و سندی—بود. و چطور به لطف ویاگرا می‌توانست همه‌ی آن‌ها را راضی نگه دارد (اگرچه پولش احتمالاً آن‌ها را به اندازه‌ی کافی راضی می‌کرد). اگر او می‌خواست با کس دیگری بخوابد، قاعده این بود که همه‌ی آن‌ها با هم این کار را انجام دهند. از این مکالمه دریافتم که این مردی بود که تمام عمر هر چقدر سکس می‌خواست داشته و در هفتاد و سه سالگی هنوز هم به دنبال زنان است. کی این متوقف می‌شود؟ اگر هیو هفنر هنوز تسلیم نشده است، من چه وقت رها خواهم شد؟

اگر «لِی‌گاید» هیچ‌وقت به راهم نمی‌افتاد، من مثل اکثر مردان هیچ وقت در تفکرم نسبت به جنس مخالف تغییر نمی‌کردم. در واقع، احتمالاً از اکثر مردان بدتر شروع کردم. در سال‌های قبل از نوجوانی، بازی دکتری نبود، هیچ دختری نبود که یک دلار برای نگاه کردن به زیر دامنشان می‌خواستند، هیچ کلاس هم‌فزونی‌هایی را در جاهایی که نباید به آن‌ها دست زده می‌شد، قلقلک نمی‌دادم. بیشتر دوران نوجوانی‌ام را محبوس در خانه بودم، بنابراین وقتی تنها فرصت جنسی نوجوانی‌ام فراهم شد—دختر مستی از سال اول دانشکده زنگ زد و پیشنهاد یک بلو جامپ به من داد—مجبور شدم امتناع کنم، وگرنه با خشم مادرم مواجه می‌شدم. در دانشگاه شروع به پیداکردن خودم کردم: چیزهایی که به آن‌ها علاقه داشتم...